\subsection{Nonrational detectors with the original affine action factor}
The intrinsic curve need not be used to change the physical action law.
A finite meromorphic clock map gives a second version that retains the
original factor $T-x$ while allowing nonrational detector primitives.
Let $\pi:X\to\mathbb P^1$ be a nonconstant meromorphic map of degree
$m$, commuting with the real structures. Let a real arc $J_X$ map
diffeomorphically onto a real clock interval $J_0$, and let
$I_0\subset\R$ be an open interval strictly to its left. Use real
meromorphic detector differentials $E$ with poles bounded by an effective
divisor $D$ avoiding $J_X$. Their primitives on $J_X$, expressed in the
clock coordinate $T=\pi(z)$, need not have rational derivatives in $T$.
The scalar law is now exactly
\[
                   f(T)=e(T)+\log(T-x)-\log(T-y),
                  \qquad x\ne y\in I_0.
\]
Put $S_\pi=\pi(\operatorname{supp}D)\cap I_0$ and define
\[
 R_c^\pi=\sum_{s\in S_\pi}c_s\frac{\dd\pi}{\pi-s},\qquad
 B_s^\pi=\frac{\dd\pi}{(\pi-s)^2}.
\]
Here $c$ runs through the same finite charge set $\mathcal Z_2(S_\pi)$.
The notation $\dd\pi$ denotes the meromorphic differential of the
clock map, not an integration variable.

\begin{theorem}[Branched-clock detector classification]
\label{thm:v86-branched}
For the original logarithmic action law in this clock chart, regular
separation holds if and only if
\[
 R_c^\pi\notin E\quad(c\in\mathcal Z_2(S_\pi)),\qquad
 E\cap\operatorname{span}\{B_s^\pi:s\in A\}=\{0\}
                    \quad(A\subset S_\pi,\ 1\le|A|\le2).
\]
Under these tests, any $\deg D+4m+2g$ distinct clocks in $J_0$ give
an injective immersion, with a smooth local inverse and compact-class
Lipschitz bounds as above. Branch values in the action interval are
allowed. The count is sufficient.
\end{theorem}
\begin{proof}
Differentiate a scalar collision on the clock arc and pull it back to
$X$. The action contribution is $R_c^\pi$ with signed charges at the
four action values, after cancellations. At a point of ramification
index $e_p$ over an action $s$, the differential
$\dd\pi/(\pi-s)$ has residue $e_p$. Distinct action fibres are
disjoint. A nonzero charge therefore produces a pole at every point
of its fibre; membership in $E$ requires, in particular, that its action
value belong to $S_\pi$. Zero charges give identical ordered pairs.
The realization of every nonzero charge by two ordered unequal pairs
is the same as in Theorem~\ref{thm:v86-divisor}. This proves the
secant criterion in both directions. For tangents,
$B_s^\pi$ has pole order $e_p+1$ at each point over $s$, and no poles
elsewhere. Poles in distinct fibres cannot cancel, which proves the
tangent criterion, including one-fibre obstructions. Constants again
remove the primitive ambiguity.

For the clock budget a secant contribution has pole degree at most
$4m$: each noncancelled action fibre contributes at most $m$.
There is no pole over infinity, since the total charge is zero.
For a tangent, each of the two possible fibres contributes
$\sum_{p\mapsto s}(e_p+1)=m+\#\pi^{-1}(s)\le2m$.
These formulas also prove the assertions at ramified action values.
Adding $D$, the pole degree of the differentiated collision or tangent
is at most $\deg D+4m$. Its zero degree, if the differential is
nonzero, is at most $\deg D+4m+2g-2$. Rolle's theorem for the stated
number of clocks forces it to vanish. The classification and the
compact inverse argument finish the proof.
\end{proof}
For a concrete positive-genus example take a nonsingular real elliptic
curve $Y^2=(T-a)(T-b)(T-c)$ with $a<b<c$ and clock interval strictly
to the right of $c$, on the positive real sheet. The clock map $\pi=T$
has degree two. The real holomorphic differential $\dd T/Y$ spans a
detector derivative space with $D=0$. Its primitive is an elliptic
integral and has a nonrational derivative in the clock coordinate.
The two obstruction families are absent, so constants and this elliptic
primitive form a regularly separating detector space for the unchanged
logarithmic action law. Ten clocks are a sufficient uniform design.
This example and the theorem concern genuinely nonrational detectors,
not only a replacement of the action curve by an Abelian integral.

